\section{Relative heads of actual transitive actions}\label{sec:relative-ranks}
For $N\lhd A$ and a prime $p$, write
\[
 h_{p,A}(N)=\dim_{\mathbb F_p}N/(N^p[N,A]),\qquad
 \rho_A(N)=h_{2,A}(N),\quad \delta_A(N)=h_{3,A}(N).
\]
The distinction between a relative head and the abelianization of an
abstract isomorphic group is essential. For $B\leq C$ normal in $A$,
the quotient word maps onto the corresponding quotient word downstairs;
its kernel is an image of the head of $B$. Consequently
\begin{equation}\label{eq:relative-head-chain}
 h_{p,A}(C)\leq h_{p,A}(B)+h_{p,A/B}(C/B).
\end{equation}
If $G\leq\prod_{i=1}^t A_i$ is full on its actual orbit actions, filter
$G$ by the kernels of projection onto coordinates $i+1,\ldots,t$.
The image of the $i$th factor is a normal subgroup $N_i\lhd A_i$.
Conjugation on that factor is exactly the action of $A_i$, since the
coordinate projection of $G$ is onto. Thus
\begin{equation}\label{eq:orbit-relative-head}
 d_p(G)\leq\sum_i h_{p,A_i}(N_i).
\end{equation}
This assertion does not add the ranks of the full coordinate images.

\begin{theorem}[Binary first head]\label{thm:first-head}
For a faithful transitive $U\leq S_w$, $w\geq2$, and $N\lhd U$,
\[
 \rho_U(N)\leq3w/8,
\]
except for the whole-action pairs $(U,U)$ of the four critical actions
in Table~\ref{tab:critical}. Their heads have dimensions $1,2,2,4$.
\end{theorem}
\begin{proof}
We use strong induction on $w$. Tracey's induced-module theorem
\cite[Theorem~1.6]{Tracey2018} gives, for a submodule of an induced
binary module with fibre dimension $a$ and index $q\geq2$, at most
$aE(q,2)\leq aq/2$ module generators. Indeed, if $q$ has odd part
greater than one, the largest-prime-power branch is at most $q/3$;
if $q=2^k$, the other branch is
$\lfloor\sqrt{2/\pi}\,q/\sqrt{k}\rfloor$, which is at most $q/2$
for $k\geq3$ and equals $1,2$ for $k=1,2$.

Suppose first that $N$ is intransitive, with $q\geq2$ orbits of size
$m$. Let $K$ be the kernel of $U$ on these orbits and let $H$ be its
transitive action on one of them. The faithful embedding of $K$ in
$H^q$ and its coordinate-kernel filtration give
\[
 \rho_U(N)\leq\rho_K(N)\leq\sum_{i=1}^q\rho_H(E_i),
 \qquad E_i\lhd H.
\]
If $H$ is not critical, induction suffices. In each critical action
every proper normal subgroup has relative head at most $3m/8$:
for $C_2,V_4$ this is immediate; in $D_8$, squares and commutators
kill its central involution in each proper noncentral normal subgroup;
in $E_8$, every noncentral normal subgroup has commutator with $E_8$
equal to its centre and has image of dimension at most three modulo
that centre. These observations also suffice if the common image of
$N$ on an orbit is proper in $H$.

The remaining case is a subdirect $N\leq H^q$ with $H$ critical.
The block stabilizer induces the actual local automorphisms on every
characteristic elementary section of $H$; the sum of its conjugate
sections is its induced module of index $q$. For $C_2,V_4$ the preceding
module bound gives $q/2,q$. For $D_8,E_8$, use both
$N\cap Z(H)^q$ and its image in $(H/Z(H))^q$. Their fibre dimensions
are respectively $(1,2)$ and $(1,4)$, so
\eqref{eq:relative-head-chain} gives $3q/2,5q/2$, respectively.
All four bounds are at most $3mq/8$. In particular the central line
of a dihedral block has not been omitted.

If $N$ is transitive, then
\[
 \rho_U(N)\leq d(N)\leq
 \left\lfloor\frac{\sqrt3\,w}{2\sqrt{\log_2w}}\right\rfloor
\]
by \cite[Theorem~1.1]{Tracey2021Sharp}. This is at most $3w/8$ for
$w\geq38$, and also for
\[w\in\{16,19,24,27,30,32,33,35,36\}.\]
The remaining finite input checks the whole heads in the complete
transitive catalogues of degrees
\[
 2,3,\ldots,23,25,26,28,29,31,34,37.
\]
Its only failures are the four critical actions. The catalogue
classification is that of \cite{HoltRoyle2020,TransGrp365}; this slice
has 7,259 action classes. A separate enumeration of all 233 transitive
normal pairs in degrees $2,4,8$ shows that a critical transitive $N$
has a head above $3w/8$ only when $U=N$ is the corresponding critical
action. Thus the whole-head test is applied to the actual action of
$N$, and the separate pair test handles its possible overgroups.
The finite calculation constructs $N/(N^2[N,U])$ and verifies its
elementary binary invariants; its precise bounded contract is recorded
in Appendix~\ref{sec:certificates}. Together with the preceding
induction this proves the assertion in all degrees.
\end{proof}

We next give the relative ternary bounds needed for the one-triple
case. Let $\mathcal E_9$ consist of the four transitive actions
$9T2,9T6,9T7,9T17$, of orders $9,27,27,81$.

\begin{lemma}[Primitive inputs and block filtration]\label{lem:ternary-block}
For a primitive action $R$ of degree $r$, let $a_3(R)$ count its
composition factors of order three, with multiplicity. Then
\[
 a_3(R)\leq r/3,
\]
with equality only for the natural $C_3,S_3$ actions and the affine
degree-nine actions of orders $216,432$. Outside these four actions,
$a_3(R)<3r/10$. Every primitive normal pair has
$\delta_R(N)<3r/20$, apart from the whole natural $C_3,A_4$ pairs.
For a transitive $A$ with a minimal block system of $s$ blocks of size
$r$, primitive block component $R$, top $T$, and block kernel $K$,
\begin{equation}\label{eq:ternary-block-recurrence}
 \delta_A(N)\leq a_3(R)E(s,3)+\delta_T(NK/K),\qquad
 E(s,3)=\min\left\{
 \left\lfloor\frac{s}{\sqrt{\pi\log_3s_3}}\right\rfloor,
 \frac{s}{\operatorname{lpp}(s/s_3)}\right\}.
\end{equation}
The first entry is infinite if $s_3=1$.
\end{lemma}
\begin{proof}
The primitive generator bound
\cite[Theorem~1.1]{HoltRoneyDougal2013} gives
$\delta_R(N)\leq\lfloor\log_2r\rfloor$; its exceptional natural
$S_3$ case has zero relative ternary head. This proves the strict
$3r/20$ bound for $r\geq34$: check $34\leq r\leq63$, and then
use $20j<3\cdot2^j$ for $j\geq6$. The composition-length bound
\cite[Theorem~1.3]{GlasbyPraegerRosaVerret2018} gives
$a_3(R)\leq(8/3)\log_2r-4/3<3r/10$ for $r\geq45$.
At $45$ this follows from $45^{16}<2^{89}$, and the difference
is increasing thereafter. The bounded inputs enumerate all primitive
normal pairs through degree 33 and all primitive composition series
through degree 44. These are respectively 945 normal pairs in 253
actions and 336 actions, with exactly the exceptions stated above.
The weaker $r/3$ assertion follows immediately, including those
exceptions. Classification completeness uses
\cite{RoneyDougal2005Primitive,PrimGrp344}; the finite computations
count order-three composition factors, rather than the ternary
valuation of group order.

For the block assertion, choose a chief series $R_i$ of $R$ and
intersect $A$ and then $N$ with the subgroups $R_i^s$ of the wreath
base. Tracey's local-chief construction
\cite[Lemma~5.8]{Tracey2018} makes each elementary chief layer a
submodule of the actual module induced from a block stabilizer;
a nonabelian layer is zero or a perfect chief factor. A ternary
layer of fibre dimension $a_i$ contributes at most $a_iE(s,3)$
by the induced-module theorem. Other primes and perfect layers
contribute zero. Sum this single chain using
\eqref{eq:relative-head-chain}, and then its top quotient. The
ternary fibre dimensions sum to $a_3(R)$, proving
\eqref{eq:ternary-block-recurrence}.
\end{proof}

\begin{theorem}[Ternary relative stability]\label{thm:rank3-stability}
Except for the whole natural $C_3,A_4$ pairs and the four whole
$\mathcal E_9$ pairs, every transitive normal pair of degree $w$ obeys
\[
 \delta_A(N)\leq\lfloor5w/27\rfloor.
\]
Consequently the bounds $w/4$ and $2w/9$ hold outside, respectively,
the whole natural $C_3$ pair and the two whole natural pairs.
Positive equality in the former occurs only at natural $A_4$;
equality in the latter occurs precisely at the four whole
$\mathcal E_9$ pairs.
For an arbitrary $G\leq S_m$ with $c$ actual regular $C_3$ orbits
and $a$ actual natural $A_4$ orbits,
\begin{equation}\label{eq:rank3-pattern}
 9d_3(G)\leq2m+3c+a.
\end{equation}
\end{theorem}
\begin{proof}
The finite primitive input in Lemma~\ref{lem:ternary-block} already
implies the assertion in the primitive case. The finite transitive
normal-pair input in degrees $9,18$ says that the only two-ninth
equalities are the four whole $\mathcal E_9$ pairs, that every other
degree-nine pair has head at most one, and that every degree-eighteen
pair has head at most three. These are complete finite catalogues,
not assumptions on the possible block systems.

Directly from the formula for $E$ one has
\[
 E(s,3)\leq10s/27\quad(s\notin\{2,6,18\}),\qquad
 E(s,3)\leq7s/18\quad(s\notin\{2,6\}),
\]
and $E(2),E(3),E(4),E(6),E(9),E(18)=1,1,1,3,3,7$.
For example, if the prime-to-three part is neither one nor two,
the largest-prime-power branch is at most $s/4$; otherwise write
$s=3^k$ or $2\cdot3^k$, and for $k\geq3$ use $\pi k>9$.
The finitely many remaining powers give the displayed values. In
particular
\[
 E(s,3)+\lfloor5s/27\rfloor\leq5s/9
 \quad(s\notin\{2,6\}).
\]
Use induction in \eqref{eq:ternary-block-recurrence}. For an ordinary
top and $s\notin\{2,6\}$, a block of size two has no ternary layer;
a block of size three uses the last inequality; for $r\geq4$ use
\[
 9\delta_A(N)\leq7rs/6+5s/3\leq5rs/3.
\]
For a whole $\mathcal E_9$ top, the bound is
$\delta\leq3a_3(R)+2\leq r+2\leq5r/3$ for $r\geq3$;
$r=2$ is immediate. For $s=2$ the bound is $a_3(R)\leq r/3$.
For $s=6$ it is $3a_3(R)+1$: use the degree-eighteen input at
$r=3$, the strict composition-density list for $4\leq r\leq8$,
and $r+1\leq10r/9$ for $r\geq9$. Finally natural $C_3$ and
$A_4$ tops give $a_3(R)+1$. In the $C_3$ case use the degree-nine
input at $r=3$, $a_3(R)\leq1$ at $r=4$, and
$r/3+1\leq5r/9$ for $r\geq5$; the $A_4$ case is immediate.
These cases exhaust the induction.

The coarser conclusions now follow by comparing $5/27<2/9<1/4$
and inspecting the six stated exceptions. Apply the two-ninth bound
to \eqref{eq:orbit-relative-head}; a natural $C_3$ layer costs three
and a natural $A_4$ layer one above twice its degree. This proves
\eqref{eq:rank3-pattern}. In particular a no-$C_3$, no-$A_4$
residual has $d_3(G)\leq2m/9$; neither exclusion can be omitted.
\end{proof}

\begin{lemma}[Two small module refinements]\label{lem:ternary-refinements}
In \eqref{eq:ternary-block-recurrence}, if the top contains a
semiregular $C_3$, the kernel term can be replaced by $a_3(R)s/3$.
If the top contains a soluble transitive subgroup, it can be replaced
by $a_3(R)\beta(s)$, where
\[
 \beta(s)=[x^{\lfloor k(s)/2\rfloor}]
       \prod_{q\mid s}(1+x+\cdots+x^{q-1})^{v_q(s)},\qquad
 k(s)=\sum_{q\mid s}v_q(s)(q-1).
\]
\end{lemma}
\begin{proof}
Lift a generator of the semiregular $C_3$ to a $3$-element $x$ of
the original group. Its order can exceed three. On an induced ternary
module of fibre dimension $a$, bases of one fibre in each three-cycle
generate the module over $\mathbb F_3[\langle x\rangle]$ with
$as/3$ vectors. This algebra is the chain ring
$\mathbb F_3[t]/(t^{3^e})$. For its finite modules, the number of
generators is the dimension of the socle; a submodule has smaller
socle, and a quotient cannot require more generators. This bounds
the actual subquotient in the chief filtration, including fibre twists.

For the second assertion retain the product-of-chains transversal in
the proof of \cite[Theorem~4.13, Lemma~4.15(ii),
Proposition~4.18]{Tracey2018}. Its rank polynomial is the product
displayed above. The chain product has a symmetric chain
decomposition, so its width is its middle coefficient $\beta(s)$.
The leading-height argument in that proof chooses at most $a$ module
generators per member of an antichain of this transversal. Keeping
its actual width, before the subsequent binomial estimate, gives
$a\beta(s)$. This uses the soluble-transitive-subgroup hypothesis
on the actual coset action.
\end{proof}

\begin{theorem}[Necessary high ternary actions]\label{thm:three-twentieths}
If a transitive normal pair $(A,N)$ of degree $w$ satisfies
$20\delta_A(N)>3w$, its actual action is one of the following:
\begin{enumerate}
\item natural $C_3$ or natural $A_4$;
\item $6T1,6T4,6T5,6T6$;
\item $12T20,12T85,12T164,12T194,12T228,12T229,12T265$;
\item a transitive $3$-group of degree $9$ or $27$.
\end{enumerate}
This is a necessary list of actions, not a claim that all their normal
pairs have high relative head.
\end{theorem}
\begin{proof}
Primitive actions are covered by Lemma~\ref{lem:ternary-block}.
Use a minimal block system and put $v=a_3(R)$ in
\eqref{eq:ternary-block-recurrence}. For a nonexceptional top and
$s\notin\{2,6,18\}$, $E(s,3)\leq s/3$. If $r\geq5$, then
\[
 \frac{vs/3+5s/27}{rs}\leq\frac19+\frac5{27r}
 \leq\frac4{27}<\frac3{20};
\]
at $r=4$, $v\leq1$ gives $14s/27<3(4s)/20$.
At $s=18$ use $7v+3$; this is at most $27r/10$ for $r\geq9$,
and $v\leq1,1,1,2,2$ handles $r=4,5,6,7,8$.
Every transitive degree-six group contains a semiregular $C_3$:
its Sylow $3$-subgroup either has order three and trivial point
stabilizers, or is $C_3^2$ on two triples. Thus at $s=6$,
Lemma~\ref{lem:ternary-refinements} gives $2v+1\leq9r/10$ for
$r\geq4$. At $s=2$ the top head vanishes and the stronger primitive
density bound gives $v\leq3r/10$, except possibly $r=9$.
The complete degree-eighteen normal-pair input gives head at most two
there and removes that exception.

For a natural $C_3$ top the bound is $v+1$: the only $r\geq4$
case not below $9r/20$ is $r=4$, giving degree twelve. A natural
$A_4$ top gives $v+1\leq3r/5$. A whole $\mathcal E_9$ top gives
$3v+2\leq27r/20$ for $r\geq6$, and $v\leq1$ deals with $r=4,5$.
For binary blocks $v=0$, only the natural $C_3$ top can fail,
giving degree six.

It remains to treat ternary blocks, $r=3$. For a nonexceptional top,
the scalar bound $E(s,3)+\lfloor5s/27\rfloor\leq\lfloor9s/20\rfloor$
leaves, beyond $s=2,3,4,6,9$, only $18,27,54,81,162$.
To verify completeness, write $s=3^kt$. If $t\notin\{1,2\}$,
use $1/4+5/27<9/20$. If $t\in\{1,2\}$ and $k\geq5$,
use $1/\sqrt{15}<143/540=9/20-5/27$. The remaining powers give
exactly the stated finite set. The chain widths at its five degrees
are $5,7,13,19,35$; adding $\lfloor5s/27\rfloor$ gives
$8,12,23,34,65$, all below the required bounds. This handles tops
containing a soluble transitive subgroup. In prime-power degree a
Sylow subgroup is transitive (compare its order with a point
stabilizer), so it handles $s=27,81$ automatically.

For completeness the nonsoluble seams use the following bounded
inputs: every nonsoluble degree-eighteen action has a semiregular
$C_3$ and all its relative heads are at most one; and the maxima of
$a_3(R)$ for primitive degrees $6,9,18,27,54$ are at most $0,3,0,7,0$.
The former checks all 91 nonsoluble actions and their 521 normal
pairs, with an actual semiregular permutation witness for each
action; the latter checks the 38 primitive actions at those degrees.
They imply the following structural bounds for nonsoluble transitive
groups, without a degree-54 or degree-162 transitive census:
\[
 \max_{N\lhd A}\delta_A(N)\leq7\quad(\deg A=54),\qquad
 \max_{N\lhd A}\delta_A(N)\leq24\quad(\deg A=162).
\]
Indeed the possible minimal block sizes at degree 54 are
$2,3,6,9,18,27$, and their kernel-plus-top bounds are respectively
$5,7,2,7,1,7$. At degree 162 the sizes are
$2,3,6,9,18,27,54,81$, with bounds
$15,24,5,24,2,22,1,15$. For the size-three rows nonsolubility
forces the top to be nonsoluble, because the local component is
$C_3$ or $S_3$. The last row uses
$(8/3)\log_2(81)-4/3<16$ and integrality. Primitive cases have
the smaller bounds $\lfloor\log_2(54)\rfloor=5$ and
$\lfloor\log_2(162)\rfloor=7$.
The ternary child bounds at top degrees $18,54,162$ are therefore
$6+1=7$, $17+7=24$, $45+24=69$, respectively, each sufficient.

Only child degrees $6,9,12,18,27$ remain. The complete normal-pair
catalogues at degrees $6,12,18$ contain 1,300 actions and 20,410
pairs; their strict high-action list is exactly items two and three,
with no degree-eighteen hit. Degree nine is already covered by
Theorem~\ref{thm:rank3-stability}. At degree 27 a top head at most
one gives $E(9,3)+1=4\leq81/20$. Otherwise the top is a whole
$\mathcal E_9$ action, hence a $3$-group. Embed the action in
$S_3\wr T$ and put $V=C_3^9$. If the image of the full block kernel
in $C_2^9$ is nontrivial, its coordinate characters on $V$ are all
nontrivial, by top transitivity. Maschke's theorem then makes every
invariant translation subquotient have zero trivial head under
that binary image. The remaining kernel quotient is binary, so
its ternary relative head is zero, and the whole head is at most two.
If that image is trivial, the full block kernel and the top are both
$3$-groups, making the actual degree-27 action a $3$-group. This
exhausts the final case and proves the theorem.
\end{proof}
